% auto-generated by experiments/make_report.py -- do not edit
\newcommand{\cicIsolationForestacc}{77.22}
\newcommand{\cicIsolationForestaccsd}{0.21}
\newcommand{\cicIsolationForestprec}{74.00}
\newcommand{\cicIsolationForestprecsd}{0.25}
\newcommand{\cicIsolationForestrec}{95.51}
\newcommand{\cicIsolationForestrecsd}{0.27}
\newcommand{\cicIsolationForestfOne}{83.39}
\newcommand{\cicIsolationForestfOnesd}{0.13}
\newcommand{\cicIsolationForestfpr}{50.09}
\newcommand{\cicIsolationForestfprsd}{0.73}
\newcommand{\cicIsolationForestauc}{80.64}
\newcommand{\cicIsolationForestaucsd}{0.86}
\newcommand{\cicIsolationForesttime}{0.0}
\newcommand{\cicRFacc}{99.48}
\newcommand{\cicRFaccsd}{0.08}
\newcommand{\cicRFprec}{99.59}
\newcommand{\cicRFprecsd}{0.03}
\newcommand{\cicRFrec}{99.55}
\newcommand{\cicRFrecsd}{0.11}
\newcommand{\cicRFfOne}{99.57}
\newcommand{\cicRFfOnesd}{0.07}
\newcommand{\cicRFfpr}{0.62}
\newcommand{\cicRFfprsd}{0.04}
\newcommand{\cicRFauc}{99.88}
\newcommand{\cicRFaucsd}{0.03}
\newcommand{\cicRFmfOne}{97.73}
\newcommand{\cicRFmfOnesd}{0.10}
\newcommand{\cicRFtime}{0.3}
\newcommand{\cicXGBoostacc}{99.62}
\newcommand{\cicXGBoostaccsd}{0.08}
\newcommand{\cicXGBoostprec}{99.63}
\newcommand{\cicXGBoostprecsd}{0.08}
\newcommand{\cicXGBoostrec}{99.74}
\newcommand{\cicXGBoostrecsd}{0.06}
\newcommand{\cicXGBoostfOne}{99.68}
\newcommand{\cicXGBoostfOnesd}{0.06}
\newcommand{\cicXGBoostfpr}{0.56}
\newcommand{\cicXGBoostfprsd}{0.13}
\newcommand{\cicXGBoostauc}{99.96}
\newcommand{\cicXGBoostaucsd}{0.01}
\newcommand{\cicXGBoostmfOne}{97.35}
\newcommand{\cicXGBoostmfOnesd}{0.98}
\newcommand{\cicXGBoosttime}{0.2}
\newcommand{\cicMLPacc}{97.74}
\newcommand{\cicMLPaccsd}{0.22}
\newcommand{\cicMLPprec}{97.15}
\newcommand{\cicMLPprecsd}{0.38}
\newcommand{\cicMLPrec}{99.13}
\newcommand{\cicMLPrecsd}{0.23}
\newcommand{\cicMLPfOne}{98.13}
\newcommand{\cicMLPfOnesd}{0.18}
\newcommand{\cicMLPfpr}{4.34}
\newcommand{\cicMLPfprsd}{0.60}
\newcommand{\cicMLPauc}{99.77}
\newcommand{\cicMLPaucsd}{0.01}
\newcommand{\cicMLPmfOne}{93.29}
\newcommand{\cicMLPmfOnesd}{0.43}
\newcommand{\cicMLPtime}{0.1}
\newcommand{\cicCNNLSTMacc}{97.64}
\newcommand{\cicCNNLSTMaccsd}{0.19}
\newcommand{\cicCNNLSTMprec}{97.45}
\newcommand{\cicCNNLSTMprecsd}{0.25}
\newcommand{\cicCNNLSTMrec}{98.64}
\newcommand{\cicCNNLSTMrecsd}{0.32}
\newcommand{\cicCNNLSTMfOne}{98.04}
\newcommand{\cicCNNLSTMfOnesd}{0.16}
\newcommand{\cicCNNLSTMfpr}{3.86}
\newcommand{\cicCNNLSTMfprsd}{0.40}
\newcommand{\cicCNNLSTMauc}{99.57}
\newcommand{\cicCNNLSTMaucsd}{0.05}
\newcommand{\cicCNNLSTMmfOne}{89.33}
\newcommand{\cicCNNLSTMmfOnesd}{1.58}
\newcommand{\cicCNNLSTMtime}{2.1}
\newcommand{\cicLSTMacc}{94.76}
\newcommand{\cicLSTMaccsd}{0.65}
\newcommand{\cicLSTMprec}{94.86}
\newcommand{\cicLSTMprecsd}{0.24}
\newcommand{\cicLSTMrec}{96.47}
\newcommand{\cicLSTMrecsd}{0.90}
\newcommand{\cicLSTMfOne}{95.66}
\newcommand{\cicLSTMfOnesd}{0.55}
\newcommand{\cicLSTMfpr}{7.80}
\newcommand{\cicLSTMfprsd}{0.33}
\newcommand{\cicLSTMauc}{98.88}
\newcommand{\cicLSTMaucsd}{0.14}
\newcommand{\cicLSTMmfOne}{85.40}
\newcommand{\cicLSTMmfOnesd}{3.19}
\newcommand{\cicLSTMtime}{3.7}
\newcommand{\cicTransformeracc}{98.97}
\newcommand{\cicTransformeraccsd}{0.27}
\newcommand{\cicTransformerprec}{98.82}
\newcommand{\cicTransformerprecsd}{0.34}
\newcommand{\cicTransformerrec}{99.48}
\newcommand{\cicTransformerrecsd}{0.12}
\newcommand{\cicTransformerfOne}{99.15}
\newcommand{\cicTransformerfOnesd}{0.23}
\newcommand{\cicTransformerfpr}{1.78}
\newcommand{\cicTransformerfprsd}{0.51}
\newcommand{\cicTransformerauc}{99.86}
\newcommand{\cicTransformeraucsd}{0.07}
\newcommand{\cicTransformermfOne}{95.29}
\newcommand{\cicTransformermfOnesd}{0.82}
\newcommand{\cicTransformertime}{10.6}
\newcommand{\cicHALXacc}{99.50}
\newcommand{\cicHALXaccsd}{0.09}
\newcommand{\cicHALXprec}{99.48}
\newcommand{\cicHALXprecsd}{0.07}
\newcommand{\cicHALXrec}{99.68}
\newcommand{\cicHALXrecsd}{0.08}
\newcommand{\cicHALXfOne}{99.58}
\newcommand{\cicHALXfOnesd}{0.07}
\newcommand{\cicHALXfpr}{0.77}
\newcommand{\cicHALXfprsd}{0.11}
\newcommand{\cicHALXauc}{99.95}
\newcommand{\cicHALXaucsd}{0.01}
\newcommand{\cicHALXmfOne}{95.86}
\newcommand{\cicHALXmfOnesd}{0.44}
\newcommand{\cicHALXtime}{11.0}
\newcommand{\cicablfullfOne}{99.54}
\newcommand{\cicablfullfpr}{0.84}
\newcommand{\cicablfullmfOne}{95.98}
\newcommand{\cicablnoattentionfOne}{99.58}
\newcommand{\cicablnoattentionfpr}{0.75}
\newcommand{\cicablnoattentionmfOne}{96.89}
\newcommand{\cicablnoxgboostfOne}{97.07}
\newcommand{\cicablnoxgboostfpr}{6.86}
\newcommand{\cicablnoxgboostmfOne}{90.37}
\newcommand{\cicablnorawfeaturesfOne}{99.37}
\newcommand{\cicablnorawfeaturesfpr}{1.16}
\newcommand{\cicablnorawfeaturesmfOne}{94.01}
\newcommand{\cicablnosmotefOne}{99.59}
\newcommand{\cicablnosmotefpr}{0.78}
\newcommand{\cicablnosmotemfOne}{96.78}
\newcommand{\cicablnoprivacyfOne}{99.72}
\newcommand{\cicablnoprivacyfpr}{0.65}
\newcommand{\cicablnoprivacymfOne}{94.84}
\newcommand{\cicabllstmonlyfOne}{97.07}
\newcommand{\cicabllstmonlyfpr}{5.87}
\newcommand{\cicabllstmonlymfOne}{87.96}
\newcommand{\unswIsolationForestacc}{73.24}
\newcommand{\unswIsolationForestaccsd}{0.69}
\newcommand{\unswIsolationForestprec}{69.27}
\newcommand{\unswIsolationForestprecsd}{1.58}
\newcommand{\unswIsolationForestrec}{92.73}
\newcommand{\unswIsolationForestrecsd}{3.45}
\newcommand{\unswIsolationForestfOne}{79.23}
\newcommand{\unswIsolationForestfOnesd}{0.19}
\newcommand{\unswIsolationForestfpr}{50.64}
\newcommand{\unswIsolationForestfprsd}{5.75}
\newcommand{\unswIsolationForestauc}{81.53}
\newcommand{\unswIsolationForestaucsd}{0.47}
\newcommand{\unswIsolationForesttime}{0.1}
\newcommand{\unswRFacc}{86.67}
\newcommand{\unswRFaccsd}{0.09}
\newcommand{\unswRFprec}{81.18}
\newcommand{\unswRFprecsd}{0.10}
\newcommand{\unswRFrec}{98.66}
\newcommand{\unswRFrecsd}{0.06}
\newcommand{\unswRFfOne}{89.07}
\newcommand{\unswRFfOnesd}{0.07}
\newcommand{\unswRFfpr}{28.03}
\newcommand{\unswRFfprsd}{0.20}
\newcommand{\unswRFauc}{97.86}
\newcommand{\unswRFaucsd}{0.02}
\newcommand{\unswRFmfOne}{49.04}
\newcommand{\unswRFmfOnesd}{0.31}
\newcommand{\unswRFtime}{1.3}
\newcommand{\unswXGBoostacc}{86.81}
\newcommand{\unswXGBoostaccsd}{0.14}
\newcommand{\unswXGBoostprec}{81.57}
\newcommand{\unswXGBoostprecsd}{0.18}
\newcommand{\unswXGBoostrec}{98.22}
\newcommand{\unswXGBoostrecsd}{0.02}
\newcommand{\unswXGBoostfOne}{89.13}
\newcommand{\unswXGBoostfOnesd}{0.10}
\newcommand{\unswXGBoostfpr}{27.18}
\newcommand{\unswXGBoostfprsd}{0.33}
\newcommand{\unswXGBoostauc}{98.26}
\newcommand{\unswXGBoostaucsd}{0.03}
\newcommand{\unswXGBoostmfOne}{49.93}
\newcommand{\unswXGBoostmfOnesd}{0.84}
\newcommand{\unswXGBoosttime}{1.2}
\newcommand{\unswMLPacc}{86.36}
\newcommand{\unswMLPaccsd}{0.24}
\newcommand{\unswMLPprec}{81.39}
\newcommand{\unswMLPprecsd}{1.00}
\newcommand{\unswMLPrec}{97.57}
\newcommand{\unswMLPrecsd}{1.32}
\newcommand{\unswMLPfOne}{88.73}
\newcommand{\unswMLPfOnesd}{0.05}
\newcommand{\unswMLPfpr}{27.38}
\newcommand{\unswMLPfprsd}{2.15}
\newcommand{\unswMLPauc}{97.75}
\newcommand{\unswMLPaucsd}{0.49}
\newcommand{\unswMLPmfOne}{43.70}
\newcommand{\unswMLPmfOnesd}{1.38}
\newcommand{\unswMLPtime}{0.2}
\newcommand{\unswCNNLSTMacc}{86.13}
\newcommand{\unswCNNLSTMaccsd}{0.35}
\newcommand{\unswCNNLSTMprec}{80.63}
\newcommand{\unswCNNLSTMprecsd}{0.42}
\newcommand{\unswCNNLSTMrec}{98.47}
\newcommand{\unswCNNLSTMrecsd}{0.38}
\newcommand{\unswCNNLSTMfOne}{88.66}
\newcommand{\unswCNNLSTMfOnesd}{0.26}
\newcommand{\unswCNNLSTMfpr}{28.98}
\newcommand{\unswCNNLSTMfprsd}{0.82}
\newcommand{\unswCNNLSTMauc}{98.17}
\newcommand{\unswCNNLSTMaucsd}{0.10}
\newcommand{\unswCNNLSTMmfOne}{42.43}
\newcommand{\unswCNNLSTMmfOnesd}{0.35}
\newcommand{\unswCNNLSTMtime}{2.1}
\newcommand{\unswLSTMacc}{84.99}
\newcommand{\unswLSTMaccsd}{0.05}
\newcommand{\unswLSTMprec}{79.25}
\newcommand{\unswLSTMprecsd}{0.24}
\newcommand{\unswLSTMrec}{98.54}
\newcommand{\unswLSTMrecsd}{0.39}
\newcommand{\unswLSTMfOne}{87.85}
\newcommand{\unswLSTMfOnesd}{0.01}
\newcommand{\unswLSTMfpr}{31.61}
\newcommand{\unswLSTMfprsd}{0.58}
\newcommand{\unswLSTMauc}{97.94}
\newcommand{\unswLSTMaucsd}{0.13}
\newcommand{\unswLSTMmfOne}{40.76}
\newcommand{\unswLSTMmfOnesd}{0.88}
\newcommand{\unswLSTMtime}{4.5}
\newcommand{\unswTransformeracc}{86.94}
\newcommand{\unswTransformeraccsd}{2.60}
\newcommand{\unswTransformerprec}{82.28}
\newcommand{\unswTransformerprecsd}{3.89}
\newcommand{\unswTransformerrec}{97.66}
\newcommand{\unswTransformerrecsd}{1.48}
\newcommand{\unswTransformerfOne}{89.23}
\newcommand{\unswTransformerfOnesd}{1.78}
\newcommand{\unswTransformerfpr}{26.18}
\newcommand{\unswTransformerfprsd}{7.31}
\newcommand{\unswTransformerauc}{98.09}
\newcommand{\unswTransformeraucsd}{0.29}
\newcommand{\unswTransformermfOne}{43.69}
\newcommand{\unswTransformermfOnesd}{1.49}
\newcommand{\unswTransformertime}{12.2}
\newcommand{\unswHALXacc}{86.78}
\newcommand{\unswHALXaccsd}{0.08}
\newcommand{\unswHALXprec}{81.56}
\newcommand{\unswHALXprecsd}{0.07}
\newcommand{\unswHALXrec}{98.20}
\newcommand{\unswHALXrecsd}{0.13}
\newcommand{\unswHALXfOne}{89.11}
\newcommand{\unswHALXfOnesd}{0.07}
\newcommand{\unswHALXfpr}{27.21}
\newcommand{\unswHALXfprsd}{0.14}
\newcommand{\unswHALXauc}{98.25}
\newcommand{\unswHALXaucsd}{0.05}
\newcommand{\unswHALXmfOne}{48.27}
\newcommand{\unswHALXmfOnesd}{0.50}
\newcommand{\unswHALXtime}{13.9}
\newcommand{\unswablfullfOne}{89.20}
\newcommand{\unswablfullfpr}{27.22}
\newcommand{\unswablfullmfOne}{47.70}
\newcommand{\unswablnoattentionfOne}{88.89}
\newcommand{\unswablnoattentionfpr}{27.38}
\newcommand{\unswablnoattentionmfOne}{48.44}
\newcommand{\unswablnoxgboostfOne}{88.96}
\newcommand{\unswablnoxgboostfpr}{27.34}
\newcommand{\unswablnoxgboostmfOne}{41.87}
\newcommand{\unswablnorawfeaturesfOne}{88.93}
\newcommand{\unswablnorawfeaturesfpr}{27.13}
\newcommand{\unswablnorawfeaturesmfOne}{45.01}
\newcommand{\unswablnosmotefOne}{89.22}
\newcommand{\unswablnosmotefpr}{26.62}
\newcommand{\unswablnosmotemfOne}{48.09}
\newcommand{\unswablnoprivacyfOne}{89.17}
\newcommand{\unswablnoprivacyfpr}{26.89}
\newcommand{\unswablnoprivacymfOne}{49.15}
\newcommand{\unswabllstmonlyfOne}{87.50}
\newcommand{\unswabllstmonlyfpr}{33.13}
\newcommand{\unswabllstmonlymfOne}{39.82}
\newcommand{\unswthru}{4{,}885}
\newcommand{\unswlatmed}{3.74}
\newcommand{\unswlatpn}{4.88}
\newcommand{\unswprivus}{0.73}
\newcommand{\unswparams}{61{,}690}
\newcommand{\cicthru}{3{,}899}
\newcommand{\ciclatmed}{3.19}
\newcommand{\ciclatpn}{4.81}
\newcommand{\cicprivus}{0.62}
\newcommand{\cicparams}{61{,}255}
\newcommand{\cicfpsharelowmedium}{69.2}
\newcommand{\cicalertsharehigh}{98.4}
\newcommand{\cictopOneZeroattackmass}{41.9}
\newcommand{\cicsevlow}{56.7}
\newcommand{\cicsevnlow}{30}
\newcommand{\cicsevmedium}{90.4}
\newcommand{\cicsevnmedium}{52}
\newcommand{\cicsevhigh}{99.8}
\newcommand{\cicsevnhigh}{5{,}020}
\newcommand{\cicsevbenign}{0.4}
\newcommand{\cicsevnbenign}{3{,}398}
\newcommand{\unswfpsharelowmedium}{87.3}
\newcommand{\unswalertsharehigh}{77.1}
\newcommand{\unswtopOneZeroattackmass}{57.8}
\newcommand{\unswsevlow}{22.8}
\newcommand{\unswsevnlow}{6{,}813}
\newcommand{\unswsevmedium}{38.2}
\newcommand{\unswsevnmedium}{5{,}702}
\newcommand{\unswsevhigh}{97.0}
\newcommand{\unswsevnhigh}{42{,}156}
\newcommand{\unswsevbenign}{2.6}
\newcommand{\unswsevnbenign}{27{,}661}
\newcommand{\PRIVCOSTMAX}{0.11}
\newcommand{\cicgapfOne}{3.9}
\newcommand{\cicfprratio}{10}
\newcommand{\cicprivcost}{0.11}
\newcommand{\cicIsolationForesttprOne}{1.1}
\newcommand{\cicRFtprOne}{99.6}
\newcommand{\cicXGBoosttprOne}{99.8}
\newcommand{\cicMLPtprOne}{94.6}
\newcommand{\cicCNNLSTMtprOne}{94.2}
\newcommand{\cicLSTMtprOne}{80.4}
\newcommand{\cicTransformertprOne}{99.0}
\newcommand{\cicHALXtprOne}{99.8}
\newcommand{\unswgapfOne}{1.3}
\newcommand{\unswfprratio}{1}
\newcommand{\unswprivcost}{-0.03}
\newcommand{\unswIsolationForesttprOne}{0.5}
\newcommand{\unswRFtprOne}{77.7}
\newcommand{\unswXGBoosttprOne}{85.9}
\newcommand{\unswMLPtprOne}{81.0}
\newcommand{\unswCNNLSTMtprOne}{84.9}
\newcommand{\unswLSTMtprOne}{81.9}
\newcommand{\unswTransformertprOne}{82.6}
\newcommand{\unswHALXtprOne}{86.9}


\begin{table*}[t]
\caption{Binary detection (attack vs.\ benign, $\tau=0.5$) and multi-class macro-F1 (M-F1) on the test
sets, mean{\tiny$\pm$std} over 3 seeds, in \%. Best mean per column and dataset in bold.
All methods receive identical privatized inputs ($k=10$, $\varepsilon=1$).}
\label{tab:main}
\centering
\footnotesize
\setlength{\tabcolsep}{4.2pt}
\begin{tabular}{@{}l@{\hspace{3pt}}ccccccc@{}}
\toprule
\textbf{Method} & \textbf{Acc.} & \textbf{Prec.} & \textbf{Rec.} & \textbf{F1} & \textbf{FPR} & \textbf{AUC} & \textbf{M-F1} \\
\midrule
\multicolumn{8}{l}{\textit{CICIDS2017 (stratified subset)}} \\
Isolation Forest~\cite{liu2008iforest} & 77.38 & 74.05 & 95.80 & 83.53 & 50.12 & 81.20 & -- \\
Random Forest~\cite{breiman2001rf} & 99.41 & \textbf{99.57} & 99.45 & 99.51 & \textbf{0.65} & 99.85 & 97.63 \\
XGBoost~\cite{chen2016xgboost} & \textbf{99.53} & 99.51 & \textbf{99.71} & \textbf{99.61} & 0.73 & \textbf{99.94} & \textbf{98.02} \\
MLP & 97.65 & 96.79 & 99.37 & 98.06 & 4.93 & 99.78 & 92.79 \\
CNN-LSTM~\cite{halbouni2022cnnlstm} & 97.65 & 97.92 & 98.15 & 98.04 & 3.11 & 99.57 & 90.54 \\
\bottomrule
\end{tabular}

\end{table*}

\subsection{Detection Performance}
Table~\ref{tab:main} compares HALX with the seven baselines.

\textbf{CICIDS2017.} Every supervised tree-based method is close to the
ceiling. HALX reaches F1 \cicHALXfOne\% (FPR \cicHALXfpr\%), statistically
level with Random Forest (\cicRFfOne\%) and within 0.1 points of
XGBoost (\cicXGBoostfOne\%). The pure deep models fall clearly behind:
the Transformer reaches \cicTransformerfOne\%, CNN-LSTM \cicCNNLSTMfOne\% and the attention-free
BiLSTM \cicLSTMfOne\%, with 2--10$\times$ higher false-positive rates. Inside
HALX the XGBoost head lifts the BiLSTM encoder by
\UNSWorCICgap{} F1 points and cuts its FPR by an order of magnitude. This
confirms that on flow features the trees should make the final
decision~\cite{grinsztajn2022tree}.

\textbf{UNSW-NB15.} On the official split all methods lose about ten
accuracy points. The test set has a different benign/attack mix and
benign traffic that is not represented in training, so at $\tau=0.5$
every supervised model has a high FPR
(\unswXGBoostfpr\% for XGBoost, \unswHALXfpr\% for HALX). This
well-documented gap~\cite{catillo2022transferability,engelen2021troubleshooting}
is far more realistic than the near-perfect CICIDS2017 numbers, and
\UNSWsummary{}

\begin{figure}[t]
\centering
% ROC curves (log-scaled FPR), seed 0, both datasets.
\begin{tikzpicture}
\begin{groupplot}[
  group style={group size=2 by 1, horizontal sep=7mm, y descriptions at=edge left},
  halxbase, width=0.62\columnwidth, height=0.44\columnwidth,
  xmode=log, xmin=0.01, xmax=100, ymin=50, ymax=100.5,
  xlabel={FPR (\%, log scale)}, ylabel={TPR (\%)},
  xtick={0.01,0.1,1,10,100}, xticklabels={0.01,0.1,1,10,100},
  legend columns=4, legend to name=rocleg,
]
\nextgroupplot[title={CICIDS2017}]
\addplot[sTR] table[x=fpr,y=tpr] {generated/roc_cic_Transformer.dat};
\addplot[sRF] table[x=fpr,y=tpr] {generated/roc_cic_RF.dat};
\addplot[sXGB] table[x=fpr,y=tpr] {generated/roc_cic_XGBoost.dat};
\addplot[sHALX] table[x=fpr,y=tpr] {generated/roc_cic_HALX.dat};
\legend{Transformer, Random Forest, XGBoost, HALX}
\nextgroupplot[title={UNSW-NB15}]
\addplot[sTR] table[x=fpr,y=tpr] {generated/roc_unsw_Transformer.dat};
\addplot[sRF] table[x=fpr,y=tpr] {generated/roc_unsw_RF.dat};
\addplot[sXGB] table[x=fpr,y=tpr] {generated/roc_unsw_XGBoost.dat};
\addplot[sHALX] table[x=fpr,y=tpr] {generated/roc_unsw_HALX.dat};
\end{groupplot}
\node[anchor=north] at ($(group c1r1.south)!0.5!(group c2r1.south)+(0,-7mm)$) {\pgfplotslegendfromname{rocleg}};
\end{tikzpicture}

\caption{ROC curves on the test sets (seed 0, log-scaled FPR). The low-FPR
region corresponds to the alert budget of a SOC.}
\label{fig:roc}
\end{figure}

\textbf{Operating point.} Fig.~\ref{fig:roc} shows that the ranking quality
(AUC) of the tree-based models and HALX is much better than their
$\tau=0.5$ FPR suggests. At a 1\% false-alarm budget on UNSW-NB15, HALX
detects \UNSWtprHALX\% of attacks, compared with \UNSWtprXGB\% for XGBoost
and \UNSWtprRF\% for Random Forest (seed 0). In deployment, $\tau$ should therefore be
set to the SOC's alert budget, not left at 0.5.

\subsection{Ablation Study}
\begin{table}[t]
\caption{Ablation of HALX components (test F1, FPR and macro-F1, \%).
CICIDS2017: 3 seeds; UNSW-NB15: \ABLunswSeeds{}.}
\label{tab:ablation}
\centering
\scriptsize
\setlength{\tabcolsep}{2.6pt}
\begin{tabular}{@{}l@{\hspace{2pt}}cccccc@{}}
\toprule
 & \multicolumn{3}{c}{\textbf{CICIDS2017}} & \multicolumn{3}{c}{\textbf{UNSW-NB15}} \\
\cmidrule(lr){2-4}\cmidrule(l){5-7}
\textbf{Configuration} & F1 & FPR & M-F1 & F1 & FPR & M-F1 \\
\midrule
Full HALX & -- & -- & -- & 89.20 & 27.22 & 47.70 \\
w/o attention (mean pooling) & -- & -- & -- & -- & -- & -- \\
w/o XGBoost (dense softmax head) & -- & -- & -- & -- & -- & -- \\
w/o raw features in head & -- & -- & -- & -- & -- & -- \\
w/o SMOTE & -- & -- & -- & -- & -- & -- \\
w/o privacy layer & -- & -- & -- & -- & -- & -- \\
BiLSTM only & -- & -- & -- & -- & -- & -- \\
\bottomrule
\end{tabular}

\end{table}
\ABLATIONTEXT{}

\subsection{Privacy--Utility Trade-off}
\begin{figure}[t]
\centering
% Privacy-utility trade-off: F1 and FPR vs epsilon (k=10) and vs k (eps=1).
\begin{tikzpicture}
\begin{groupplot}[
  group style={group size=2 by 1, horizontal sep=6mm},
  halxbase, width=0.6\columnwidth, height=0.42\columnwidth,
  xtick={0,1,2,3,4},
  every axis plot/.append style={mark size=1.6pt},
]
\nextgroupplot[title={DP budget ($k{=}10$)}, xlabel={$\varepsilon$}, ylabel={$\Delta$F1 (pp)},
  legend style={at={(0.97,0.04)}, anchor=south east, fill=white, draw=none, row sep=-2pt},
  xticklabels={0.01,0.1,1,10,$\infty$}]
\addplot[sXGB, mark=square*] table[x=x,y=df1] {generated/priv_unsw_XGBoost_eps.dat};
\addplot[sRF, mark=triangle*] table[x=x,y=df1] {generated/priv_cic_XGBoost_eps.dat};
\IfFileExists{generated/priv_cic_HALX_eps.dat}{\addplot[sHALX, mark=*] table[x=x,y=df1] {generated/priv_cic_HALX_eps.dat};}{}
\legend{XGBoost (UNSW), XGBoost (CIC), HALX (CIC)}
\nextgroupplot[title={Anonymity level ($\varepsilon{=}1$)}, xlabel={$k$},
  xticklabels={1,5,25,100,500}]
\addplot[sXGB, mark=square*] table[x=x,y=df1] {generated/priv_unsw_XGBoost_k.dat};
\addplot[sRF, mark=triangle*] table[x=x,y=df1] {generated/priv_cic_XGBoost_k.dat};
\IfFileExists{generated/priv_cic_HALX_k.dat}{\addplot[sHALX, mark=*] table[x=x,y=df1] {generated/priv_cic_HALX_k.dat};}{}
\end{groupplot}
\end{tikzpicture}

\caption{Test F1 vs.\ DP budget $\varepsilon$ (left, $k=10$) and
anonymity level $k$ (right, $\varepsilon=1$). $\varepsilon=\infty$ and
$k=1$ disable the respective mechanism.}
\label{fig:privacy}
\end{figure}
\PRIVACYTEXT{}

\subsection{Explainability and Severity Triage}
\begin{figure}[t]
\centering
\input{figures/fig_attention_unsw}
\caption{Mean attention-pooling weight of the ten features most attended
for attack flows (UNSW-NB15 test set, seed 0), compared with benign flows.
A uniform weight would be $1/42\approx2.4\%$.}
\label{fig:attention}
\end{figure}
\EXPLAINTEXT{}

\subsection{Computational Cost}
\COSTTEXT{}
